\documentclass[11pt,a4paper]{article}
\usepackage[utf8]{inputenc}
\usepackage[margin=1in]{geometry}
\usepackage{amsmath, amssymb, amsfonts, amsthm, bm}
\usepackage{booktabs}
\usepackage{hyperref}
\usepackage{microtype}
\usepackage{array}

\newtheorem{proposition}{Proposition}
\newtheorem{lemma}{Lemma}
\newtheorem{remark}{Remark}

\newcommand{\E}{\mathbb{E}}
\newcommand{\R}{\mathbb{R}}

\title{\textbf{Why Successor-Measure Methods Fail at Manipulation\\[2pt]
\large An investigation of FB, PSM and CRL on cube manipulation, and the interaction-weighted fix}}
\date{\today}

\begin{document}
\maketitle

\begin{abstract}
\noindent
Forward--Backward representations (FB), the Proto Successor Measure (PSM) and Contrastive RL (CRL)
all estimate the same object --- the discounted successor measure --- and all underperform a plain
goal-conditioned TD baseline (GCIQL) on cube manipulation by 20--45 success points. We trace this
gap through a cascade of diagnostics down to three precise defects: (i) the FB objective allocates
representation capacity by \emph{temporal persistence} under the data, which deletes the
interaction variable (cube held / free / slipping) that manipulation hinges on; (ii) a value
trained against \emph{independently sampled} goals receives no signal about controllability,
regardless of readout architecture; (iii) the task vector $z$ acquires two inconsistent meanings
--- a reward descriptor for the value and a policy index for the actor. All three defects live in
the \emph{training distribution}, not in the estimators: every estimator-side correction we tried
(better task inference, better readout architectures on frozen features) recovers nothing. We
derive two distribution-level fixes --- an interaction-weighted goal measure (IWR, ported from
contrastive learning) and trajectory-future task sampling (zHER, which is exactly CRL's coupling)
--- and show they target complementary defects.
\end{abstract}

\section{Setup: one object, four estimators}\label{sec:setup}

An MDP with states $s$, actions $a$, discount $\gamma$; an offline dataset of transitions
$(s,a,s')$ from a behavior policy $\beta$ (OGBench \texttt{cube-single} play data: a scripted arm
repeatedly pick-and-placing a cube). $\rho$ is the dataset distribution, $\rho_G$ the marginal of
next-states (the ``goal pool'').

The central object is the \textbf{successor measure} of a policy $\pi$,
\begin{equation}
M^{\pi}(s,a,\,\cdot\,) \;=\; \sum_{k\ge 1}\gamma^{\,k-1}\,P\!\left(s_{t+k}\in\cdot \mid s_t=s,\,a_t=a,\,\pi\right),
\label{eq:sm}
\end{equation}
the discounted distribution over future states. It determines everything: for any reward $r$,
$Q^{\pi}_r(s,a)=\int M^{\pi}(s,a,dg)\,r(g)$. ``Zero-shot'' methods learn a representation of $M$
once and read off $Q_r$ for any $r$ at test time. We write $m^{\pi}(s,a,g)$ for the density of
$M^{\pi}$ with respect to $\rho_G$.

\subsection{FB: Forward--Backward representations}
FB factorizes the density \emph{linearly},
\begin{equation}
m^{\pi_z}(s,a,g)\;\approx\;F(s,a,z)^{\top}B(g),\qquad F,B\in\R^{d},
\end{equation}
jointly over a policy family $\{\pi_z\}_{z\in\R^d}$. Training minimizes the measure-matching TD
loss
\begin{equation}
\mathcal{L}_{\rho}(F,B)=
\underbrace{\E_{(s,a)\sim\rho}\,\E_{g\sim\rho_G}\!\Big[\big(F(s,a,z)^{\top}B(g)-\gamma\,\bar F(s',\pi_z(s'),z)^{\top}\bar B(g)\big)^2\Big]}_{\text{off-diagonal: where error is charged}}
\;-\;2\,\underbrace{\E_{(s,a,s')\sim\rho}\big[F(s,a,z)^{\top}B(s')\big]}_{\text{diagonal (Dirac term)}}
\label{eq:fbloss}
\end{equation}
subject to the orthonormalization $\E_{\rho_G}[BB^{\top}]=I$ (coefficient
\texttt{ortho\_coef}), with $z$ drawn 50/50 from a random sphere and from $z=B(g)$, $g\sim\rho_G$.
At test time a task with reward $r$ is encoded and evaluated by
\begin{equation}
z_r=\E_{\rho_G}[\,r(g)\,B(g)\,], \qquad Q_r(s,a)=F(s,a,z_r)^{\top}z_r,
\label{eq:fbinfer}
\end{equation}
and the agent executes $\pi_{z_r}$. Both identities are exact \emph{iff} $m=F^{\top}B$ exactly
and $B$ is orthonormal under $\rho_G$.

\subsection{PSM: Proto Successor Measure}
PSM learns the same factorization in two stages: a basis $\varphi(g)\in\R^{d}$ trained by a
measure loss of the form \eqref{eq:fbloss} against a family of \emph{proto-policies} (random
behaviors), with the same orthonormalization on $\varphi$; then successor features
$\psi(s,a,z)\approx\E_{\pi_z}\big[\sum_{k}\gamma^{k-1}\varphi(s_{t+k})\big]$ trained by TD with
$\varphi$ detached. The value is again bilinear, $Q(s,a,z)=\psi(s,a,z)^{\top}z$; inference mirrors
\eqref{eq:fbinfer} with $\varphi$ for $B$. The actor (FlowBC: a flow behavior-cloning policy with
a $Q$-ascent term) never reshapes $\varphi$.

\subsection{CRL: Contrastive RL}
CRL learns a critic $f(s,a,g)=\phi(s,a)^{\top}\psi(g)$ by InfoNCE, where \emph{positives} are
trajectory futures $g=s_{t+k}$, $k\sim\mathrm{Geom}(1-\gamma)$, and negatives are random states.
The NCE optimum is the \emph{log}-density ratio:
\begin{equation}
f^{*}(s,a,g)\;=\;\log\frac{p(g\mid s,a)}{\rho_G(g)}+\mathrm{const}
\;=\;\log m^{\beta}(s,a,g)+\mathrm{const},
\end{equation}
using that the positive distribution is $p(g\mid s,a)=(1-\gamma)\,m^{\beta}(s,a,g)\,\rho_G(g)$
(see \eqref{eq:coupling} below). So CRL is a \emph{log-linear} model of the same object
\eqref{eq:sm}, trained on a \emph{different coupling} of states and goals.

\subsection{GCIQL: the reference that works}
Goal-conditioned IQL: a monolithic MLP value $Q(s,a,g)$, expectile TD, goals sampled HER-style
($62.5\%$ trajectory futures). No factorization, no zero-shot transfer --- but it sets the
ceiling.

\paragraph{Structural differences at a glance.}
\begin{center}\small
\begin{tabular}{@{}lllll@{}}
\toprule
 & represents & form & state--goal coupling & task inference\\
\midrule
FB & $m$ (density) & bilinear $F^{\top}B$ & independent ($\rho\otimes\rho_G$) & $z_r=\E[rB]$\\
PSM & $m$ on basis $\varphi$ & bilinear $\psi^{\top}z$ & independent & $z_r=\E[r\varphi]$\\
CRL & $\log m$ & bilinear $\phi^{\top}\psi$ & \textbf{trajectory futures} & none (goal given)\\
GCIQL & $Q_g$ directly & \textbf{MLP} & \textbf{trajectory futures} & none (goal given)\\
\bottomrule
\end{tabular}
\end{center}
The investigation shows the \emph{coupling} column, not the \emph{form} column, is decisive.

\section{The failure}\label{sec:failure}

All methods: 10 seeds, 1M steps, state observations, identical data. IQM success over the five
cube tasks:
\begin{center}\small
\begin{tabular}{@{}lccccc@{}}
\toprule
 & GCIQL & FB & PSM & CRL & RLDP\\
\midrule
IQM success (\%) & $68.3\pm3.7$ & $50.0\pm7.0$ & $36.7\pm4.5$ & $25.9\pm2.0$ & $22.7\pm6.0$\\
\bottomrule
\end{tabular}
\end{center}
A phase diagnostic decomposes each rollout into \emph{reach $\to$ acquire control $\to$ maintain
control $\to$ transport}. Every method reaches the cube ($\approx100\%$) and most acquire a grasp
($80$--$90\%$); failures concentrate \emph{after} control acquisition (PSM: $35\%$ of episodes end
having lost the grasp, $6\%$ holding but misplaced). Manipulation is not hard because of approach
or grasping; the missing capability is \emph{keeping and using} control. Training longer does not
help (FB $50.0\!\to\!47.7$, PSM $36.7\!\to\!40.3$ at 1.5M): the ceiling is structural.

\section{Localizing the defect: a diagnostic cascade}\label{sec:cascade}

Each step eliminates one hypothesis. All numbers are 10-seed means on the same checkpoints.

\subsection{The basis is not missing the information}\label{sec:basis}
A \emph{factorization probe} regresses task variables from the learned goal embedding: the
cube-to-goal distance is recoverable from PSM's $\varphi$ with $R^{2}=0.90$ (FB's $B$: $0.21$).
A \emph{$z$-decode probe} asks what the training task vectors $z=B(g)$ encode about their goal
state $g$ (held-out $R^{2}$ / AUC):
\begin{center}\small
\begin{tabular}{@{}lcccc@{}}
\toprule
$z$ encodes $\to$ & cube position & cube lift & hand--cube dist. & gripper contact\\
\midrule
FB ($d{=}50$) & $0.30$ & $\mathbf{0.02}$ & $0.09$ & $0.75$\\
PSM ($d{=}128$, ortho $10^{3}$) & $0.90$ & $\mathbf{0.81}$ & $0.76$ & $0.97$\\
\bottomrule
\end{tabular}
\end{center}
FB's task vocabulary cannot express ``the cube is lifted'': its tasks are coarse positions. PSM's
tasks encode the full manipulation configuration --- and PSM \emph{still} fails. A good basis is
necessary but not sufficient.

\subsection{The value is blind to controllability}\label{sec:blind}
A \emph{counterfactual value probe}: reset the simulator to $200$ in-hand grasps from the data,
roll out each agent's own policy, and ask whether the agent's value ranks grasps by their actual
success probability (rank-AUC): GCIQL $0.89$; CRL $0.68$; FB $\approx0.60$; RLDP $0.52$; PSM
$0.49$. A matched-pair \emph{aliasing probe} (nearest fail-bound vs.\ success-bound look-alike
states) agrees: bilinear values separate them at $0.61$, near chance. The quantity the policy
maximizes does not know whether a grasp will be kept --- even when (PSM) its own basis does.

\subsection{The actor is exonerated}
Replacing GCIQL's TD3 actor with the \emph{same} flow-BC actor stack used by FB/PSM yields IQM
$72.3$ --- above GCIQL-TD3 ($68.3$). The behavior-cloning prior supplies grasping competence to
every method; the differentiator is purely the steering value. Actor architecture is not the
bottleneck.

\subsection{Eval-time inference cannot fix it}\label{sec:ladder}
If $m\neq F^{\top}B$ exactly, the rule \eqref{eq:fbinfer} is no longer optimal. We tested
corrected rules on frozen checkpoints ($10$ seeds $\times$ $5$ tasks, paired): \textbf{L1}, the
generalized-least-squares projection $z=(\E[BB^{\top}])^{-1}\E[rB]$ (the exact $L_2$ projection of
$r$ onto $\mathrm{span}(B)$ when $B$ is not orthonormal); \textbf{L2}, a \emph{Bellman-consistent}
$z$ minimizing the model's own Bellman residual over the single vector $z$.
\begin{center}\small
\begin{tabular}{@{}lcccc@{}}
\toprule
success (\%) & L0 (standard) & L1 (GLS) & L2 (Bellman) & L2 (balanced)\\
\midrule
FB & $45.1$ & $+2.7$ (n.s.) & $-33.4^{*}$ & $-33.8^{*}$\\
PSM (ortho $10^{3}$) & $39.1$ & $-0.3$ & $-26.8^{*}$ & $-22.9^{*}$\\
PSM (ortho $1$) & $27.2$ & $-22.0^{*}$ & $-13.2^{*}$ & $-13.5^{*}$\\
\bottomrule
\end{tabular}
\end{center}
Success is monotone in the cosine to the standard $z$: \emph{any} departure degrades performance,
including the mathematically more correct ones. Interpretation: at eval, $z$ is not a reward
descriptor but an \emph{index into the policy family}, and the actor only understands indices from
its training distribution. The value's task embedding and the policy's task index have
\textbf{decoupled semantics} --- a training-level defect no inference rule can repair. (An earlier
goal-\emph{measure} fix, reweighting toward rare cube positions, failed the same way: it raised
the best achievable ranking direction $z^{*}$ but left $z=B(g)$ at chance.)

\subsection{The readout form is not the constraint}\label{sec:readout}
The natural suspect for \S\ref{sec:blind} is the bilinear form (a rank-$d$ inner product) versus
GCIQL's MLP. Control: train an MLP value $Q(s,a,z)$ on the \emph{frozen} PSM basis, same
conditioning $z$, same data, but with FB's goal sampling (independent uniform goals + a Dirac
fraction). Result: counterfactual AUC $0.48/0.49/0.52$ across seeds --- statistically identical to
the bilinear ($0.48/0.53/0.42$), flat from 25k steps on. \textbf{An MLP given FB's training signal
converges to the same blindness.} The signal, not the function class, is deficient. GCIQL's
$0.89$ then differs in exactly one remaining ingredient: its goals are trajectory futures.

\subsection{The objective is the active culprit (the ortho paradox)}\label{sec:ortho}
Empirically, the best configuration for \emph{both} FB and PSM sets \texttt{ortho\_coef}
$=10^{3}$ --- which looks like turning the successor-measure loss off. A gradient audit (per-term
gradients on $B/\varphi$ across checkpoints; the loss terms separate exactly by linearity of the
gradient) refines this:
\begin{center}\small
\begin{tabular}{@{}lll@{}}
\toprule
 & $C\|\nabla\mathcal{L}_{\mathrm{ortho}}\|\,/\,\|\nabla\mathcal{L}_{\mathrm{measure}}\|$ & consequence for the basis\\
\midrule
FB, $C{=}10^{3}$ & $1.2$--$1.9$, flat over training & measure term fully active --- and $B$ is lift-blind\\
PSM, $C{=}10^{3}$ & $26\to4.5$ (whitening locks in early) & basis $\approx$ isometry; interaction preserved ($0.81$)\\
PSM, $C{=}1$ & $0.01$ & measure term in control; basis degrades ($0.65$); success drops\\
\bottomrule
\end{tabular}
\end{center}
For FB the measure loss is \emph{not} suppressed: it shapes $B$ throughout training and still
produces an interaction-blind basis. For PSM, whenever the measure loss controls the basis, the
basis gets \emph{worse}. The objective itself prefers the wrong structure. Section
\ref{sec:mech} explains why.

\section{Mechanism}\label{sec:mech}

\subsection{Persistence ordering: what a rank-$d$ successor model keeps}\label{sec:persistence}
At its optimum, the loss \eqref{eq:fbloss} makes $F^{\top}B$ the best rank-$d$ approximation of
$m$ in $L_2(\rho\otimes\rho_G)$: the top-$d$ singular space of the successor operator. Writing
$M=\sum_{k\ge1}\gamma^{k-1}P^{k}$, a state variable that mixes quickly under the behavior process
(eigenvalue $\lambda_P$ of $P$ small) contributes $\sim(1-\gamma\lambda_P)^{-1}$ --- a small
singular value --- and each singular direction is further weighted by its data mass. Hence
\begin{equation}
\text{capacity ordering}\;\propto\;\underbrace{\text{temporal persistence under }\beta}_{\text{slow variables win}}
\;\times\;\underbrace{\rho\text{-mass}}_{\text{common states win}}.
\label{eq:ordering}
\end{equation}
In play data the arm/cube \emph{positions} drift slowly and are always present; the interaction
variable (free/held/slipping) toggles every few dozen steps, and its boundary --- the contact set
$\mathcal{I}$ --- has little mass. By \eqref{eq:ordering} it is truncated first. This is the
measured lift-$R^{2}=0.02$ of \S\ref{sec:basis}: \textbf{the objective allocates capacity by
persistence, which in manipulation is anti-correlated with control relevance.} It also resolves
the ortho paradox of \S\ref{sec:ortho}: strong whitening is persistence-agnostic --- it pushes the
basis toward an isometry that keeps \emph{all} directions, including interaction --- so
suppressing the objective's preference helps.

\subsection{The coupling: why independent goals erase controllability}\label{sec:coupling}
Let $q(g\mid s,a)$ be the goal distribution used in training, and consider two anchors that alias
in observation space: $s_{\mathrm{hold}}$ (stable grasp) and $s_{\mathrm{slip}}$ (about to drop).
The pressure on \emph{any} value regression to distinguish them is
\begin{equation}
D_q\;=\;\E_{g\sim q}\Big[\big(m(s_{\mathrm{hold}},g)-m(s_{\mathrm{slip}},g)\big)^{2}\Big].
\label{eq:dq}
\end{equation}
Under $q=\rho_G$ (independent goals --- FB, PSM, and the \S\ref{sec:readout} control), almost all
goals are unreachable from \emph{either} anchor within the horizon: both densities vanish,
$D_q\approx0$, the regression target is \emph{constant along the controllability direction}, and
any estimator --- bilinear or MLP --- faithfully learns to ignore it. Under trajectory-future
goals, $q$ concentrates exactly on goals reachable from $s_{\mathrm{hold}}$ but not
$s_{\mathrm{slip}}$ (transported-cube futures), making $D_q$ large.

Formally: sampling $g=s_{t+k}$ with $k\sim\mathrm{Geom}(1-\gamma)$, i.e.
$P(k)=(1-\gamma)\gamma^{k-1}$, induces
\begin{equation}
p(g\mid s,a)\;=\;(1-\gamma)\sum_{k\ge1}\gamma^{k-1}P\big(s_{t+k}=g\mid s,a\big)
\;=\;(1-\gamma)\,m^{\beta}(s,a,g)\,\rho_G(g),
\label{eq:coupling}
\end{equation}
by definition \eqref{eq:sm}: \emph{trajectory-future goals are samples from the normalized
successor measure of the behavior policy} --- precisely CRL's positive distribution, and (up to
the HER mixture) GCIQL's. The coupling column of \S\ref{sec:setup} is now derived, and it predicts
the observed ordering: GCIQL (MLP $+$ trajectory coupling) $0.89$ $>$ CRL (restrictive log-bilinear
form, right coupling) $0.68$ $>$ FB/PSM (any form, independent coupling) $\approx0.5$.

\subsection{The deep structure: a renewal decomposition}\label{sec:renewal}
Why is manipulation special? Its dynamics are \emph{mode-switched}: the cube moves only through
contact. Let $\tau$ be the first hitting time of the contact set $\mathcal{I}$ from $(s,a)$. The
strong Markov property gives, exactly,
\begin{equation}
M^{\pi}(s,a,\cdot)\;=\;M^{\pi}_{\mathrm{pre}}(s,a,\cdot)\;+\;\E\big[\gamma^{\tau}\,M^{\pi}(s_{\tau},\cdot)\big],
\label{eq:renewal}
\end{equation}
where $M_{\mathrm{pre}}$ is the successor measure of the process killed at $\mathcal{I}$. The
cross-mode part of $M$ therefore factors \emph{through the low-dimensional interaction manifold}:
a first-passage kernel onto $\mathcal{I}$ composed with the post-contact successor measure. Each
factor is smooth within its mode; their \emph{composition} is what a single flat bilinear
$F^{\top}B$ must spend rank to imitate, and what an $L_2(\rho)$ truncation smooths away first ---
producing exactly the aliasing of \S\ref{sec:blind}. (The published switched-affine analysis of
CRL in manipulation is the local linearization of \eqref{eq:renewal}.)

\section{The fix: re-shaping the training distribution}\label{sec:fix}

The cascade leaves only training-distribution interventions. The loss \eqref{eq:fbloss} exposes
three legs where the goal measure enters --- the off-diagonal integral, the orthonormalization,
and the task distribution $z=B(g)$ --- and each defect maps to a leg. All fixes are implemented
behind default-off flags (baselines byte-identical), with unit tests.

\subsection{IWR: interaction-weighted goal measure (basis leg)}
Re-weight goals by contact proximity, porting Interaction-Weighted Resampling from contrastive
learning. With effector / cube positions $u_p(g),y_p(g)$ and $d(g)=\|u_p(g)-y_p(g)\|_2$,
\begin{equation}
w(g)\;\propto\;\varepsilon+\exp\!\Big(-\tfrac{(d(g)-c)^{2}}{2\sigma^{2}}\Big),\qquad \E[w]=1,
\label{eq:iwr}
\end{equation}
with $c=0.03$ (the measured hand--cube distance of grasped states), $\sigma=0.05$, and
$\varepsilon=0.1$ keeping full support. Changing $\rho_G\to\rho'$ with density $w$ is exactly FB in
$L_2(\rho')$: by \eqref{eq:ordering} it raises the mass factor of modes supported near
$\mathcal{I}$, re-ordering the spectrum so interaction modes survive truncation. The weight flows
through the off-diagonal term, the orthonormalization, and optionally (\texttt{weight\_z}) the
task distribution --- the latter required by \S\ref{sec:ladder}: a better basis is useless if the
actor never practices interaction tasks. (The earlier \emph{coverage} reweight failed because it
tilted toward rare cube \emph{positions} --- the subspace FB already prefers. The tilt must target
the interaction boundary.)

\subsection{zHER: trajectory-future task sampling (coupling leg)}
Replace the independent goals in the task leg by each transition's own discounted future:
$z=B(s_{t+k})$, $k\sim\mathrm{Geom}(1-\gamma)$, \emph{aligned} with its anchor (not permuted). By
\eqref{eq:coupling} this gives FB exactly CRL's coupling on its task distribution: the actor
practices tasks feasible from its own states; the $z$-conditioned value trains where $D_q$
\eqref{eq:dq} is large; and eval-time task vectors --- which encode \emph{achievable} tasks ---
fall inside the training cloud, repairing the \S\ref{sec:ladder} decoupling at its source. The
measure loss, orthonormalization, inference rule, and the random-sphere half of the $z$ mix are
untouched, so the zero-shot machinery is preserved.

\subsection{Pre-registered predictions (dissociation) --- and the IWR verdict}
The fixes target complementary defects and should dissociate: IWR repairs the \emph{basis
spectrum} (predict: lift-$R^{2}$ of $B$ and interaction separability rise; modest success gain)
but not the coupling; zHER repairs the \emph{value and task legs} (predict: counterfactual AUC and
success rise) but leaves the basis persistence-ordered (lift-$R^{2}$ stays low). If both hold, the
complete fix is their composition.

\paragraph{Measured (IWR, 3 seeds @500k).} The representation-level prediction is confirmed
exactly --- and the success prediction fails in the informative direction:
\begin{center}\small
\begin{tabular}{@{}lcccc@{}}
\toprule
 & lift-$R^{2}$ of $B$ & separability AUC & success @500k &\\
\midrule
FB plain ($d{=}50$) & $0.02$ & $0.74$ & $42.0$ &\\
FB $+$ IWR ($d{=}50$) & $\mathbf{0.79}$ & $\mathbf{0.89}$ & $22.7$ &\\
FB $+$ IWR $+$ \texttt{weight\_z} & $0.77$ & $0.89$ & $16.7$ &\\
\bottomrule
\end{tabular}
\end{center}
The $\rho'$ tilt re-orders the spectrum precisely as derived (\S\ref{sec:persistence}): at
$d{=}50$ the basis now encodes lift at PSM's level ($0.79$ vs.\ $0.81$). Success halves.
Combined with the capacity control (FB at $d{=}128$: lift $0.30$, success $35.3$ vs.\ $50.0$ at
$d{=}50$; PSM $d{=}128$: lift $0.81$, success $36.7$), the completed cells form a monotone
\emph{anti}-correlation between basis interaction-encoding and success. The reading is the
thesis in its strongest form: \textbf{no in-family consumer exists for the interaction
information.} The value cannot read it (\S\ref{sec:blind}, \S\ref{sec:readout}); the inference
rule projects tasks through a basis now whitened under $\rho'$ while eval relabeling samples
$\rho$ (a GLS mismatch, \S\ref{sec:ladder}); and the policy index drifts accordingly. Improving
the representation in isolation only perturbs the one thing that was working (the BC-steered
position policy). The binding constraint is the \emph{end-to-end coupling} of representation,
value and task semantics --- which is the one ingredient the working baseline (GCIQL) has and
every retrofit lacks.

\paragraph{The mismatch is partly repairable (measured).} The GLS inference rule
$z=(\E_{\rho}[BB^{\top}])^{-1}\E_{\rho}[rB]$ applied to the \emph{IWR} checkpoints recovers
$21.0\to34.3$ ($+13.3$) and $11.0\to30.7$ ($+19.7$, \texttt{weight\_z}) --- because IWR trains
$B$ orthonormal under $\rho'$ while standard inference projects under $\rho$; the GLS solve
restores the correct inner product, and (unlike in \S\ref{sec:ladder}) the corrected $z$ stays
\emph{inside} the actor's training distribution, which was itself $\rho'$-consistent. All three
ladder regimes now behave as the theory dictates: no-op where $B$ is $\rho$-white, catastrophic
where the corrected $z$ exits the policy index cloud, strongly positive where training and
correction share the measure. The principled completion is \emph{measure-consistent inference},
$z_r=\E_{\rho'}[rB]=\E_{\rho}[w\,r\,B]$ (weight the relabel batch by the training weights):
every goal-side integral --- loss, orthonormalization, \emph{and inference} --- must use the same
measure. Even so, IWR$+$GLS ($34.3$) remains below plain FB ($42.0$): a consistent-measure,
interaction-rich basis still has no value that reads it. Completed cells of the form-$\times$-coupling table:
\begin{center}\small
\begin{tabular}{@{}lcc@{}}
\toprule
counterfactual AUC & independent goals & trajectory-future goals\\
\midrule
bilinear, linear density (FB/PSM) & $0.49$--$0.52$ (measured) & zHER --- \emph{running}\\
bilinear, log density (CRL) & --- & $0.68$ (reference)\\
MLP on frozen basis $\varphi$, relabel-$z$ & $0.50$ (measured) & $\mathbf{0.47}$ (measured) \\
MLP on raw obs, exact goal (GCIQL) & --- & $0.89$ (reference)\\
\bottomrule
\end{tabular}
\end{center}
\textbf{Update (measured):} the in-house MLP$+$trajectory cell stays at chance. Switching the
coupling alone, on a \emph{frozen basis} with \emph{relabel-inferred} task conditioning, does not
recover the signal. Two confounds separate our control from GCIQL and become the next factorial:
(i) \emph{input} --- raw observations vs.\ the frozen $\varphi(s)$; (ii) \emph{conditioning} ---
the exact task goal $g$ vs.\ the relabel-smeared $z=\E[r\varphi]$ (the deployment condition of the
FB family, which may itself cap the achievable ranking). A third candidate is end-to-end training
of the state encoder \emph{jointly} with the trajectory-coupled value --- i.e.\ the controllability
signal may only be extractable when the representation is shaped \emph{by} the coupled TD loss,
not retrofitted onto a frozen basis. Morphing the control into GCIQL one ingredient at a time
identifies which ingredient flips $0.47\to0.89$.

\section{Experimental status}\label{sec:status}
Running at matched configuration (\texttt{ortho\_coef}$=10^{3}$, $d=50$, 3 seeds, 500k):
\texttt{fb\_iwr} (pure $\rho'$), \texttt{fb\_iwr\_wz} ($+$\,\texttt{weight\_z}),
\texttt{fb\_state\_zher} (coupling), and the MLP$+$trajectory readout cell. A $2\times2$ capacity
control (FB at $d{=}128$ --- complete; PSM at $d{=}50$ --- running) separates algorithm from
embedding size in \S\ref{sec:basis}. Promotion rule: any variant beating \texttt{fb\_state\_plain}
on success \emph{or} moving its mechanism endpoint (lift-$R^{2}$, counterfactual AUC) is promoted
to $10$ seeds $\times$ 1M with the full probe suite.

\section*{Summary}
Successor-measure methods fail at manipulation for one coherent reason expressed in three places:
\emph{their training distributions never make the interaction variable matter}. The objective
ranks features by persistence and deletes it from the basis (\S\ref{sec:persistence}); independent
goal sampling makes values constant along it (\S\ref{sec:coupling}); and the task sampling never
asks the actor to use it, decoupling $z$'s two meanings (\S\ref{sec:ladder}). Estimator-side fixes
--- better inference, better readout architectures --- demonstrably recover nothing
(\S\ref{sec:ladder}, \S\ref{sec:readout}). The repairs are distributional: weight the goal measure
toward the interaction boundary (IWR), and couple tasks to trajectory futures (zHER) --- the
ingredient CRL and GCIQL already had.

\end{document}
